\documentclass[10pt,a4paper]{amsart}
\usepackage[margin=19mm]{geometry}
\usepackage{amsmath,amssymb,mathtools}
\usepackage[scr=rsfs,cal=euler]{mathalfa}
\usepackage{xfrac}
\usepackage[unicode,hidelinks,bookmarks=false]{hyperref}
\newcommand{\ie}{i.e.\ }
\newcommand{\cf}{cf.\ }
\newcommand{\resp}{respectively\ }
\newcommand{\Subsection}{\subsection}
\providecommand{\nobreakdash}{\nobreak\hbox{-}}
\setlength{\emergencystretch}{2em}
\multlinegap0pt
\usepackage[all]{xy}
\usepackage{amssymb}
\usepackage{enumerate}
\usepackage{bm}
\usepackage{algorithm}\usepackage{algpseudocode}
\usepackage{tikz}
\newtheorem{proposition}{Proposition}
\newcommand{\HH}  {\mbox{\rm H}}
\newcommand{\Nerve}{\mathcal{N}}
\title{Complexity and speed of semi-algebraic multi-persistence}
\author{Arindam Banerjee, Saugata Basu}
\date{Source: arXiv:2407.13586v2 (CC BY 4.0)}
\begin{document}
\maketitle

\noindent\emph{Source and licence.} Complexity and speed of semi-algebraic multi-persistence. Version arXiv:2407.13586v2 (CC BY 4.0), distributed by arXiv under the Creative Commons Attribution 4.0 International licence, http://creativecommons.org/licenses/by/4.0/, as recorded in the arXiv licence metadata for this exact version and shown on the abstract page https://arxiv.org/abs/2407.13586v2. Full licence notice: \texttt{licenses/arxiv-2407.13586-CC-BY-4.0.txt}.\par\smallskip

\noindent\textit{Editorial scope.} Counted unit: the article's proposition prop:refinement (source0.tex lines 1786--1790). The article's complete original proof of it is delivered with it (source0.tex lines 1792--1802, one \texttt{proof} environment of 11 lines). No mathematics is written, repaired or re-derived: the statement and the proof are the article's own lines, in the article's own notation, and no retained line is substituted or re-flowed.  Retained uncounted as the source-local context the proof needs: lines 1766--1770.  Nothing was omitted from any retained span, so the package carries no omission record.  No retained line contains a citation, so the wrapper carries no bibliography and no bibliography entry.  No retained line contains a figure environment, an image-inclusion command or drawing code, so no image is delivered and the package declares no asset dependency.  Editorial formatting only, in addition: an AMS \texttt{amsart} preamble that restates the article's own declarations and macros used by the retained lines, namely the environments \texttt{proposition} on separate counters, together with the standard packages those lines need.  The retained environments print in this document as Proposition 1 in order of appearance, whereas the article prints the same items with the numbering of its own full document.  The complete native source of the article as distributed in the arXiv e-print for this exact version is bundled unchanged as \texttt{sources/162526/source0.tex}.
\begin{proposition}[Nerve Lemma]
\label{prop:nerve}
    If $C$ is a good cover of $S$, then there is a  canonical isomorphism
    $\psi_C: \HH_*(\Nerve(C)) \rightarrow \HH_*(S)$.
\end{proposition}

\begin{proposition}
\label{prop:refinement}
   If $C' \subset C$ are two good covers of a closed and bounded semi-algebraic set $S$, then $\Nerve(C') \subset \Nerve(C)$, and the induced
   map $\HH_*(\Nerve(C')) \rightarrow \HH_*(\Nerve(C))$ is an isomorphism.
\end{proposition}

\begin{proof}
    We have the following commutative diagram 
    \[
    \xymatrix{
\HH_*(\Nerve(C')) \ar[rr]\ar[rd]^{\psi_{C'}} && \HH_*(\Nerve(C)) \ar[ld]^{\psi_C} \\
    & \HH_*(S)&
    }
  \]
where the $\psi_C',\psi_C$ are isomorphisms using Proposition~\ref{prop:nerve}.
This implies that the horizontal arrow is an isomorphism as well.
\end{proof}

\end{document}
